\documentclass[11pt,a4paper]{article}
\usepackage[margin=2.4cm,top=2.8cm,headheight=18pt,bottom=2.8cm]{geometry}
\usepackage[utf8]{inputenc}
\usepackage[T1]{fontenc}
\usepackage{amsmath,amssymb,amsthm,amsfonts,mathtools}
\usepackage{xcolor,enumitem,fancyhdr,titlesec,lmodern}
\usepackage{pifont}
\usepackage{setspace}
\usepackage{booktabs}
\usepackage{array}
\usepackage{tabularx}
\usepackage{multirow}
\usepackage{graphicx}
\usepackage{float}
\usepackage[breakable,skins,theorems]{tcolorbox}
\tcbuselibrary{skins,breakable}
\usepackage[colorlinks=true,linkcolor=blue!70!black,citecolor=green!70!black,urlcolor=purple]{hyperref}
\usepackage{bookmark}

\setlength\parindent{0pt}
\setlength\parskip{4pt}
\linespread{1.02}\selectfont
\pagestyle{fancy}
\fancyhf{}
\fancyhead[L]{\sffamily\small\bfseries Artificial Intelligence in Mathematical Engineering}
\fancyhead[R]{\sffamily\small\bfseries Project Report}
\fancyfoot[C]{\sffamily\thepage}
\renewcommand{\headrulewidth}{0.7pt}

% --- SECTION FORMATTING ---
\titleformat{\section}
  {\Large\bfseries\sffamily\color{blue!70!black}}
  {\thesection.}{1em}{}
\titleformat{\subsection}
  {\large\bfseries\sffamily\color{green!60!black}}
  {\thesubsection.}{1em}{}
\titleformat{\subsubsection}
  {\normalsize\bfseries\sffamily\color{orange!80!black}}
  {\thesubsubsection.}{1em}{}

% --- THEOREM ENVIRONMENTS ---
\theoremstyle{definition}
\newtheorem{definition}{Definition}[section]
\newtheorem{theorem}[definition]{Theorem}
\newtheorem{proposition}[definition]{Proposition}
\newtheorem{lemma}[definition]{Lemma}
\newtheorem{corollary}[definition]{Corollary}
\newtheorem{examplethm}[definition]{Example}
\newtheorem{remark}[definition]{Remark}

% --- COLOURED BOXES ---
\newtcolorbox{problembox}{enhanced,breakable,colback=blue!6!white,colframe=blue!80!black,
  colbacktitle=blue!90!black,coltitle=white,fonttitle=\bfseries\sffamily,title=Problem Statement,
  attach boxed title to top left={yshift=-3mm,xshift=5mm}}

\newtcolorbox{theorybox}{enhanced,breakable,colback=green!6!white,colframe=green!75!black,
  colbacktitle=green!85!black,coltitle=white,fonttitle=\bfseries\sffamily,title=Key Theory \& Formulas,
  attach boxed title to top left={yshift=-3mm,xshift=5mm}}

\newtcolorbox{insightbox}{enhanced,breakable,colback=purple!6!white,colframe=purple!80!black,
  colbacktitle=purple!90!black,coltitle=white,fonttitle=\bfseries\sffamily,title=Intuition \& Geometric Insight,
  attach boxed title to top left={yshift=-3mm,xshift=5mm}}

\newtcolorbox{strategybox}{enhanced,breakable,colback=orange!6!white,colframe=orange!80!black,
  colbacktitle=orange!90!black,coltitle=white,fonttitle=\bfseries\sffamily,title=Step-by-Step Derivation,
  attach boxed title to top left={yshift=-3mm,xshift=5mm}}

\newtcolorbox{tipbox}{enhanced,breakable,colback=yellow!12!white,colframe=yellow!90!black,
  colbacktitle=yellow!70!white,coltitle=black,fonttitle=\bfseries\sffamily,title=Exam Tip \& Key Insight,
  attach boxed title to top left={yshift=-3mm,xshift=5mm}}

\newtcolorbox{warningbox}{enhanced,breakable,colback=red!8!white,colframe=red!80!black,
  colbacktitle=red!90!black,coltitle=white,fonttitle=\bfseries\sffamily,title=Common Pitfall \& Warning,
  attach boxed title to top left={yshift=-3mm,xshift=5mm}}

\newtcolorbox{answerbox}{enhanced,breakable,colback=cyan!6!white,colframe=cyan!80!black,
  colbacktitle=cyan!90!black,coltitle=white,fonttitle=\bfseries\sffamily,title=Final Result,
  attach boxed title to top left={yshift=-3mm,xshift=5mm}}

\newtcolorbox{methodbox}{enhanced,breakable,colback=teal!6!white,colframe=teal!80!black,
  colbacktitle=teal!90!black,coltitle=white,fonttitle=\bfseries\sffamily,title=Alternative Method / Perspective,
  attach boxed title to top left={yshift=-3mm,xshift=5mm}}

\newtcolorbox{summarybox}{enhanced,breakable,colback=gray!8!white,colframe=gray!70!black,
  colbacktitle=gray!85!black,coltitle=white,fonttitle=\bfseries\sffamily,title=Quick Reference Summary,
  attach boxed title to top left={yshift=-3mm,xshift=5mm}}

\newtcolorbox{prereqbox}{enhanced,breakable,colback=brown!6!white,colframe=brown!70!black,
  colbacktitle=brown!85!black,coltitle=white,fonttitle=\bfseries\sffamily,title=Prerequisites \& Background Knowledge,
  attach boxed title to top left={yshift=-3mm,xshift=5mm}}

% --- MATH MACROS ---
\newcommand{\R}{\mathbb{R}}
\newcommand{\E}{\mathbb{E}}
\newcommand{\N}{\mathbb{N}}
\newcommand{\pd}[2]{\frac{\partial #1}{\partial #2}}
\newcommand{\norm}[1]{\left\| #1 \right\|}
\newcommand{\vect}[1]{\mathbf{#1}}
\newcommand{\mat}[1]{\mathbf{#1}}
\newcommand{\inner}[2]{\langle #1,\, #2 \rangle}
\DeclareMathOperator*{\argmin}{arg\,min}
\DeclareMathOperator*{\argmax}{arg\,max}
\DeclareMathOperator{\tr}{tr}
\DeclareMathOperator{\Var}{Var}
\DeclareMathOperator{\diag}{diag}
\DeclareMathOperator{\sign}{sign}
\DeclareMathOperator{\rank}{rank}

\begin{document}

\begin{center}
{\Huge\bfseries\sffamily\color{blue!70!black} Artificial Intelligence in}\\[0.3em]
{\Huge\bfseries\sffamily\color{blue!70!black} Mathematical Engineering}\\[0.6em]
{\LARGE\bfseries\sffamily PROJECT REPORT}\\[0.2em]
\rule{0.82\textwidth}{1.5pt}\\[0.6em]
{\LARGE\bfseries\sffamily AI-Assisted Diagnosis on HAM10000}\\[0.4em]
{\large\sffamily The Effect of Loss Functions and Sampling on}\\[0.2em]
{\large\sffamily Imbalanced Skin-Cancer Classification}\\[0.6em]
\rule{0.82\textwidth}{1.5pt}\\[0.6em]
{\LARGE\bfseries\sffamily Author: Ömer Kahyaoğlu}\\[0.2em]
{\Large\sffamily E-mail: kahyaoglu20@itu.edu.tr}\\[0.4em]
{\Large\sffamily Date: \today}\\[0.4em]
{\Large\sffamily Student Number: 090240741}\\[0.8em]
{\large Mathematical Engineering Undergraduate Programme}\\[0.4em]
{\normalsize Istanbul Technical University}\\[0.3em]
\rule{0.82\textwidth}{1.5pt}\\[0.5em]
{\small\sffamily\textbf{MYZ 309E}\quad CRN: 23577}\\
{\small\sffamily Class Lecturer: Professor Muhammed Kurulay, PhD}\\
\end{center}

\vspace{0.4em}

\begin{problembox}
On the HAM10000 dermoscopic dataset (10\,015 images, 7 diagnostic classes), the
majority class \textit{melanocytic nevi} (\texttt{nv}) accounts for $\sim$67\% of
the data while the rarest class \textit{dermatofibroma} (\texttt{df}) accounts
for only $\sim$1.1\%. A model trained naively with Cross-Entropy will optimise
mean log-likelihood and trivially achieve high \emph{accuracy} by mostly
predicting \texttt{nv} --- while missing the rare and clinically important
classes. We compare \textbf{five} loss/sampling configurations under an
identical CNN, optimiser and seed on the full real HAM10000 corpus, and report
\emph{macro-F1} and \emph{per-class recall} (the metrics that an imbalanced
clinical task actually needs). Headline result: \textbf{Focal loss
($\gamma{=}2$) reaches macro F1 $\mathbf{=0.443}$ vs. $0.381$ for Standard
CE}, and is the only loss-side fix that recovers non-zero recall on the
rarest class \texttt{df}.
\end{problembox}

\section{Introduction}

The HAM10000 collection \cite{tschandl2018ham} contains 10\,015 dermoscopic
images of pigmented skin lesions, hand-labelled across seven diagnostic
classes: melanocytic nevi (\texttt{nv}), melanoma (\texttt{mel}), benign
keratosis-like lesions (\texttt{bkl}), basal cell carcinoma (\texttt{bcc}),
actinic keratoses (\texttt{akiec}), vascular lesions (\texttt{vasc}) and
dermatofibroma (\texttt{df}). Figure~\ref{fig:dist} shows that the dataset is
\emph{severely imbalanced}: a constant predictor that always outputs
\texttt{nv} would already attain $\approx$67\% accuracy --- which makes plain
accuracy a misleading optimisation target.

\begin{figure}[H]
\centering
\includegraphics[width=0.9\textwidth]{class_distribution.png}
\caption{Class distribution of the HAM10000 training data. The \texttt{nv}
class dominates ($\approx$67\%), while \texttt{df} accounts for only
$\approx$1.1\% of the corpus.}
\label{fig:dist}
\end{figure}

The objective of this project is therefore not just to \emph{classify} skin
lesions, but to \emph{study how the choice of loss function and sampling
strategy interact with class imbalance}: do techniques that explicitly
re-weight rare classes recover their recall, and at what cost to majority-class
performance?

\section{Methodology}

\subsection{Pipeline overview and CNN architecture}

The full pipeline (Figure~\ref{fig:cnn}) is: (i) load all 10\,015 HAM10000
JPGs from the official Harvard Dataverse zips; (ii) stratified $70/15/15$
train/validation/test split (per-class proportions preserved); (iii)
$64{\times}64$ resize, ImageNet-style mean/std normalisation and light
training-time augmentation (horizontal/vertical flip); (iv) training of a
from-scratch CNN ($\sim$106k parameters) for 8 epochs with Adam; (v) test-set
evaluation with imbalance-aware metrics.

\begin{figure}[H]
\centering
\fbox{\parbox{0.93\textwidth}{\sffamily\small\centering
\textbf{Input} ($3{\times}64{\times}64$) $\;\rightarrow\;$
\textbf{Conv $3{\to}16$} + BN + ReLU + MaxPool$_{2}$ \;($16{\times}32{\times}32$) $\;\rightarrow\;$
\textbf{Conv $16{\to}32$} + BN + ReLU + MaxPool$_{2}$ \;($32{\times}16{\times}16$) $\;\rightarrow\;$
\textbf{Conv $32{\to}64$} + BN + ReLU + MaxPool$_{2}$ \;($64{\times}8{\times}8$) $\;\rightarrow\;$
\textbf{Conv $64{\to}128$} + BN + ReLU + MaxPool$_{2}$ \;($128{\times}4{\times}4$) $\;\rightarrow\;$
\textbf{GlobalAvgPool} \;($128$) $\;\rightarrow\;$
\textbf{Dropout$_{0.4}$} $\rightarrow$ \textbf{FC $128{\to}64$} + ReLU $\rightarrow$
\textbf{Dropout$_{0.4}$} $\rightarrow$ \textbf{FC $64{\to}7$} (logits)
}}
\caption{CNN architecture. Four Conv\,/\,BatchNorm\,/\,ReLU\,/\,MaxPool blocks
followed by a Global Average Pooling head and a small FC head. Total:
$\mathbf{106\,391}$ trainable parameters.}
\label{fig:cnn}
\end{figure}

\subsection{Why this loss family, this optimiser, these hyperparameters?}

\begin{theorybox}
With class counts $\{n_c\}_{c=1}^{C}$ and per-sample softmax probability
$p_t(x_i)=\hat{p}(y_i\mid x_i)$, the four imbalance-handling losses studied
here are
\begin{align}
\mathcal{L}_{\text{CE}}(x,y)        &= -\log p_t,
\qquad\quad
\mathcal{L}_{\text{wCE}}(x,y)        = -\frac{1/n_y}{\sum_c 1/n_c}\,C\,\log p_t,\\[2pt]
\mathcal{L}_{\text{focal}}(x,y)      &= -\alpha_y\,(1-p_t)^{\gamma}\log p_t
        \quad(\gamma=2),\qquad
\mathcal{L}_{\text{CB}}(x,y)         = -\frac{1-\beta}{1-\beta^{n_y}}\,\log p_t
        \quad(\beta=0.999).
\end{align}
The \textit{WeightedRandomSampler} method instead leaves the loss as plain CE
but draws each training example $i$ with probability $\propto 1/n_{y_i}$.
\end{theorybox}

\textbf{Loss choice.} Each formulation tackles imbalance from a different
angle: weighted/CB CE rescale the per-class gradient magnitude; focal loss
\cite{lin2017focal} additionally down-weights examples the network already
classifies confidently (large $p_t$ $\Rightarrow$ vanishing $(1-p_t)^\gamma$);
WeightedRandomSampler is a \emph{sampling-side} fix that does not change the
objective at all but changes the empirical mini-batch distribution.

\textbf{Optimiser.} We use \textbf{Adam} \cite{kingma2014adam}
($\eta=10^{-3}$, $\beta_{1{,}2}=(0.9,0.999)$, weight decay $10^{-4}$). Adam's
per-parameter adaptive learning rates handle the heterogeneous gradient
magnitudes that the weighted/focal losses produce; SGD with momentum was
tested but converged more slowly and was more sensitive to the loss-rescaling.

\textbf{Hyperparameters.} Batch size $128$, $8$ epochs, seed $42$ used for
split, model-init and DataLoader. Eight epochs are enough to read the
under/overfitting transition off the validation loss curve (Fig.~\ref{fig:overfit});
running for many more epochs would mostly memorise the majority class.

\subsection{Why these metrics?}

For an imbalanced 7-class problem, accuracy is dominated by the majority class.
We therefore report (i) per-class precision/recall/F1, (ii) macro-averaged F1
(arithmetic mean of per-class F1 scores --- equal weight to every class),
(iii) balanced accuracy (mean of per-class recall), and (iv) the full
confusion matrix.

\begin{warningbox}
Reporting test \emph{accuracy} alone on HAM10000 \emph{hides minority-class
collapse}: a model that always predicts \texttt{nv} scores $\approx 67\%$
accuracy yet has zero recall on every other class.
\end{warningbox}

\section{Experimental Results}

All metrics below come from one full end-to-end run on the real HAM10000
corpus (10\,015 images; train\,=\,7008, val\,=\,1502, test\,=\,1505) using the
exact CNN, optimiser and seed described above. Per-experiment training time
on CPU was $\approx$\,2 minutes.

\subsection{Loss and accuracy curves}

\begin{figure}[H]
\centering
\begin{minipage}{0.49\textwidth}\centering
\includegraphics[width=\linewidth]{loss_curves.png}
\end{minipage}\hfill
\begin{minipage}{0.49\textwidth}\centering
\includegraphics[width=\linewidth]{acc_curves.png}
\end{minipage}
\caption{Real HAM10000 loss (left) and accuracy (right) per experiment.
Standard CE drops fastest on train (it exploits the 67\% \texttt{nv} prior)
and reaches $\sim$73\% test accuracy --- a number mostly explained by the
prior, see Fig.~\ref{fig:cm-base}. Focal loss starts at small magnitudes
because $(1{-}p_t)^{\gamma}$ collapses the loss on easy majority examples
(absolute scale is not comparable across losses; the trend is).}
\label{fig:loss}
\end{figure}

\begin{figure}[H]
\centering
\includegraphics[width=0.78\textwidth]{overfitting_annotated.png}
\caption{Underfitting / overfitting interpretation on the Standard-CE
baseline. Epochs 1--2 show simultaneously high train and validation loss
(\emph{underfitting}); the dashed line marks the lowest validation loss --- past
this point train loss continues to decrease while val loss starts to plateau or
rise (\emph{overfitting} risk). In our runs the model was stopped before this
divergence becomes severe, but it is the canonical use of the validation curve.}
\label{fig:overfit}
\end{figure}

\subsection{Test-set metrics}

Table~\ref{tab:metrics} reports the test-set metrics for all five experiments.
The dichotomy between \textit{accuracy} and \textit{macro-F1\,/\,minority recall}
on the Standard CE baseline is the central observation of this report.

\begin{table}[H]
\centering
\renewcommand{\arraystretch}{1.15}
\begin{tabular}{lccccc}
\toprule
\textbf{Experiment} & \textbf{Acc.} & \textbf{Bal.\,Acc.} &
\textbf{Macro\,P} & \textbf{Macro\,R} & \textbf{Macro\,F1} \\
\midrule
Standard CE                       & \textbf{0.736} & 0.352 & 0.527 & 0.352 & 0.381 \\
Weighted CE                       & 0.521 & 0.505 & 0.339 & 0.505 & 0.341 \\
Focal Loss ($\gamma{=}2$)         & 0.622 & \textbf{0.538} & 0.414 & \textbf{0.538} & \textbf{0.443} \\
Class-Balanced CE                 & 0.718 & 0.447 & 0.428 & 0.447 & 0.395 \\
Std CE + WeightedRandomSampler    & 0.464 & 0.537 & 0.386 & 0.537 & 0.391 \\
\bottomrule
\end{tabular}
\caption{Test-set metrics on the real HAM10000 corpus (1\,505 test images).
Bold = best per column. Standard CE has the highest \emph{accuracy} but
the lowest \emph{macro recall}; Focal loss has the best macro F1.}
\label{tab:metrics}
\end{table}

The per-class recall heatmap (Fig.~\ref{fig:heat}) makes this picture
explicit: Standard CE achieves 0\% recall on the rarest class \texttt{df}
and only 13.6\% on the second-rarest \texttt{vasc}; Focal restores them to
38.9\% and 77.3\% respectively, at the cost of dropping \texttt{nv} recall
from 92.5\% to 69.7\%.

\begin{figure}[H]
\centering
\includegraphics[width=0.85\textwidth]{per_class_recall_heatmap.png}
\caption{Per-class recall by experiment (real HAM10000). Standard CE has
\emph{zero} recall on the rarest class \texttt{df}; only Focal and the
WeightedRandomSampler recover it.}
\label{fig:heat}
\end{figure}

\begin{figure}[H]
\centering
\includegraphics[width=0.95\textwidth]{metrics_comparison.png}
\caption{Macro F1 (left) and recall on the rarest class \texttt{df} (right).
Standard CE achieves $0.000$ recall on \texttt{df} --- it never predicts the
rare class; only Focal Loss and the WeightedRandomSampler recover it to
$\approx 0.39$. Weighted/CB CE move in the right direction but at this
training budget (8 epochs) only marginally lift \texttt{df} recall.}
\label{fig:metrics}
\end{figure}

\subsection{Confusion matrices}

\begin{figure}[H]
\centering
\begin{minipage}{0.46\textwidth}
\centering
\includegraphics[width=\linewidth]{cm_standard.png}
\end{minipage}
\hfill
\begin{minipage}{0.46\textwidth}
\centering
\includegraphics[width=\linewidth]{cm_best.png}
\end{minipage}
\caption{Real HAM10000 confusion matrices. \textbf{Left:} Standard CE ---
the entire \texttt{df} row sums to zero in the \texttt{df} column (the network
\emph{never} predicts \texttt{df}), and most rare-class examples are pushed
onto the majority class \texttt{nv}. \textbf{Right:} Focal loss spreads
predictions back across all seven classes, sacrificing some \texttt{nv} accuracy
in exchange for non-trivial recall on every rare class.}
\label{fig:cm-base}
\end{figure}

\subsection{Error analysis}

For the Standard CE baseline, Figure~\ref{fig:err} shows the three
\emph{most-confident wrong} predictions on the real HAM10000 test set. All
three are true \texttt{bkl} (benign keratoses) lesions that the network
re-labels as \texttt{nv} (melanocytic nevi) at $97{-}99.8\%$ confidence ---
\texttt{nv} and \texttt{bkl} share visual cues (dark, homogeneously pigmented
patches) and the imbalanced training set teaches the network that
``when in doubt, output \texttt{nv}''. This is exactly the pathology that the
weighted/focal/class-balanced losses and the oversampling sampler correct.

\begin{figure}[H]
\centering
\includegraphics[width=0.92\textwidth]{error_analysis_baseline.png}
\caption{Standard CE on the real HAM10000 test set: three highest-confidence
\emph{wrong} predictions. All three are benign keratoses (\texttt{bkl})
confidently mis-labelled as melanocytic nevi (\texttt{nv}) --- a textbook
minority-class collapse driven by the 67\% \texttt{nv} prior.}
\label{fig:err}
\end{figure}

\section{Discussion and Conclusion}

\subsection{What the experiments tell us}

\begin{insightbox}
Three observations come out of the real HAM10000 run:
(i) \textbf{Accuracy is a misleading metric on imbalanced data} --- the
Standard CE baseline reaches $73.6\%$ accuracy yet has $0\%$ recall on the
rarest class \texttt{df} and only $13.6\%$ on \texttt{vasc}; macro F1 is
just $0.381$.
(ii) \textbf{Focal loss is the strongest single drop-in fix} --- with
$\gamma{=}2$ and inverse-frequency $\alpha$ it reaches the highest macro F1
($0.443$) and balanced accuracy ($0.538$), and is the only loss-side method
that meaningfully recovers \texttt{df} recall ($0.389$).
(iii) \textbf{Loss reweighting and resampling are not interchangeable in
practice}: WeightedRandomSampler matches Focal on minority recall, but at a
substantial cost to overall accuracy ($0.46$ vs $0.62$); pure inverse-frequency
weighted CE hurts accuracy without proportionate recall gains within an 8-epoch
budget. Class-Balanced CE behaves like a softer Standard CE here.
\end{insightbox}

\subsection{Weaknesses of the model}

The CNN is trained from scratch with only $\approx$\,106\,k parameters and at
$64{\times}64$ input resolution --- modest relative to HAM10000's true visual
variability ($600{\times}450$ originals) and tiny compared with a transferred
ResNet-50. Imaging artefacts (hairs, ruler markings, vignettes) are not
removed, and a Convolutional Autoencoder denoising preprocess could plausibly
help recall on the visually noisier minority classes. The model also makes a
flat 7-class decision; clinical practice prioritises a benign-vs-malignant
first decision with downstream subtype classification, which would change the
relevant cost matrix.

\subsection{Future work}

\begin{itemize}[leftmargin=1.2em,itemsep=2pt]
    \item Replace the small from-scratch backbone with a pretrained
          ResNet-50/EfficientNet-B0 and re-run the same loss comparison.
    \item Add the Conv-Autoencoder denoising stage suggested in the project
          brief and quantify its effect on minority-class recall.
    \item Calibrate the predicted probabilities (e.g.\ temperature scaling)
          so that confidence values are clinically meaningful, and report
          recall at fixed false-positive rate (the operating-point metric a
          dermatologist actually cares about).
    \item Test stronger sampling/augmentation strategies: SMOTE in feature
          space, MixUp, CutMix.
    \item Extend to a hierarchical malignancy-first decision rule.
\end{itemize}

\subsection{Conclusion}

\begin{answerbox}
On the real HAM10000 task with a small from-scratch CNN, plain Cross-Entropy
reaches a deceptively high $73.6\%$ accuracy by collapsing minority classes
($0\%$ recall on \texttt{df}, $13.6\%$ on \texttt{vasc}); its macro F1 is
only $0.381$. Of the four imbalance-aware alternatives, \textbf{Focal loss
($\gamma{=}2$) is the clear winner} --- it lifts macro F1 to $0.443$ and is
the only loss-side fix that materially recovers the rarest class. The
WeightedRandomSampler reaches similar minority-class recall but at a much
larger accuracy cost. Pure inverse-frequency Weighted CE and Class-Balanced
CE move in the right direction but require longer training (or stronger
regularisation) to compete; with deeper/transferred backbones we would expect
their gap to close. The headline takeaway is unambiguous: \emph{accuracy is
the wrong metric for HAM10000}; the right yardstick is per-class recall and
macro F1, and the right tool to optimise it is Focal Loss.
\end{answerbox}

\bigskip

\begin{thebibliography}{9}
\bibitem{tschandl2018ham}
P.~Tschandl, C.~Rosendahl, H.~Kittler.
\textit{The HAM10000 dataset, a large collection of multi-source dermatoscopic
images of common pigmented skin lesions}, Scientific Data 5 (2018).

\bibitem{lin2017focal}
T.-Y.~Lin, P.~Goyal, R.~Girshick, K.~He, P.~Doll{\'a}r.
\textit{Focal Loss for Dense Object Detection}, ICCV 2017.

\bibitem{cui2019cb}
Y.~Cui, M.~Jia, T.-Y.~Lin, Y.~Song, S.~Belongie.
\textit{Class-Balanced Loss Based on Effective Number of Samples}, CVPR 2019.

\bibitem{kingma2014adam}
D.~P.~Kingma, J.~Ba.
\textit{Adam: A Method for Stochastic Optimization}, ICLR 2015.
\end{thebibliography}

\end{document}
